\section{从开环到闭环反馈}
\label{sec:16-Control-and-Reinforcement-Learning/01-Optimal-Control/01}

\subsection{淋浴困境}

昨天试出的最佳组合：冷水龙头拧三格，热水龙头拧一格。水温刚好。今天照原样拧好，打开水——要么烫得往后跳，要么凉得起鸡皮疙瘩。水厂调了供水压力，楼上邻居在你拧水那几秒冲了马桶，管道里存水的温度跟昨天已经不是同一批水。昨天的``最优刻度''，今天全盘失效。

你面前分岔出两条路。

第一条：把计划全算好——每个时刻拧多少度、拧多久——在开水之前一次性算完，合上阀门后不再管水温。计划是计划，现实是现实，井水不犯河水。

第二条：一边洗一边伸手试。凉了就拧热，烫了就拧冷。手上的感觉变成拧的动作，拧的动作又改变手上的感觉——转起来了。

这两条路的分岔口，就是开环（open-loop）与闭环（closed-loop）反馈的根本区别。判断标准很简单，跟有没有精细计算毫无关系：新的动作是否使用了刚刚观察到的结果。日常生活里``边看边调''就是反馈最朴素的形态——在你喊出``反馈''这个词之前，身体已经这样做了。

\begin{center}
\begin{tikzpicture}[>=stealth,
  block/.style={draw,rounded corners=2pt,minimum width=1.8cm,minimum height=0.75cm,fill=blue!6},
  every node/.style={font=\small}]
  \node (r) at (-4.9,0) {目标 \(r\)};
  \node[draw,circle,minimum size=0.65cm] (sum) at (-3.4,0) {\(+\)};
  \node[block] (ctrl) at (-1,0) {控制器};
  \node[block] (plant) at (2.0,0) {被控对象};
  \node (out) at (4.3,0) {输出 \(x\)};
  \draw[->,thick] (r) -- (sum);
  \draw[->,thick] (sum) -- node[above] {误差 \(e\)} (ctrl);
  \draw[->,thick] (ctrl) -- node[above] {输入 \(u\)} (plant);
  \draw[->,thick] (plant) -- (out);
  \draw[->,thick] (3.6,0) -- (3.6,-1.35) -- (-3.4,-1.35) -- (sum)
    node[pos=0.55,below] {测量反馈};
  \node at (-3.68,-0.3) {\(-\)};
  \draw[->,gray!80] (2.0,1.3) -- (plant) node[midway,right] {扰动 \(w\)};
\end{tikzpicture}

\emph{闭环的关键不是控制器多复杂，而是输出被测量并返回比较点；扰动造成的偏差因此能影响下一次动作。}
\end{center}

\subsection{先用开环把一切都算清楚}

把场景抽象成一台恒温水箱。目标 60 度。水温记 \(x(t)\)，环境温度 20 度。散热使水温按 \(-a(x - 20)\) 的速率向环境滑落，\(a = 0.1\)。加热器的功率称为输入 \(u\)：

\[
\dot x = -a(x - x_e) + u.
\]

先把所有计算做完，再合电闸。站在稳态看——水温不再变时 \(\dot x = 0\)。代入目标 60 度：

\[
0 = -0.1(60 - 20) + u \quad\Rightarrow\quad u^* = 4.
\]

维持 60 度需要功率 4。再从冷启动设计一段升温曲线——先大火烧到接近 60 度，再切到 4。整条 \(u(t)\) 在合闸前一次算完，执行期间不再用温度计。这是开环控制：输入只是时间的函数，里面不含任何关于当前水温的信息。

记偏差 \(e = x - r\)（\(r = 60\)）。把输入拆成稳态维持量加修正量：\(u = u^* + \bar u\)。代入方程，常数项恰好相消：

\[
\dot e = -a e + \bar u.
\]

开环方案取 \(\bar u \equiv 0\)——不做任何修正。若模型精确且世界安静，偏差从初始值按 \(e^{-at}\) 指数衰减到零。\(a = 0.1\) 时，每 7 秒偏差减半。纸面上一切完美。

\subsection{三道裂口：模型错了、世界推了、系统自己也在推自己}

开环的完美依赖一个前提：你对世界的描述分毫不差。现实从三个方向攻进来。

\textbf{模型测不准。}你怎么知道真实散热系数正好是 0.1？水垢厚度、环境湿度、水箱表面的灰尘——每一样都在修改散热速率。假设真实系数 \(\tilde a = 0.08\)，用 \(u^* = 4\) 去维持，真实平衡点变成

\[
x = 20 + \frac{4}{0.08} = 70\; \text{度。}
\]

你计划 60 度，实际水温蹲在 70 度不走了。而因为你从不测量水温，你永远不知道恒温箱悄悄变成了桑拿房。只要没人看一眼温度计，系统就在 70 度安静地运行到地老天荒。

\textbf{扰动不请自来。}有人开了窗，一股 \(-1\) 千瓦的冷风灌入：

\[
\dot x = -a(x - x_e) + u + w, \quad w = -1.
\]

偏差方程变为 \(\dot e = -a e + w\)。令 \(\dot e = 0\) 求新稳态：\(e = w/a = -10\) 度。水温从 60 度滑到 50 度，恒定停留在那里。你的计划依然在纸面上``正常运行''——它不可能知道窗开了，因为没有任何一条信道把``窗开了''的信息送回到计算里。

\textbf{被控对象本身可能带正反馈。}前面两个例子里散热是负反馈——偏差越大，向平衡点的恢复力越强。换一种场景：放热反应中温度越高反应越快，偏差越大放热越猛，方程变号：

\[
\dot e = +\lambda e, \quad \lambda > 0.
\]

这不是``没有恢复力''——这是有排斥力。假设 \(\lambda = 0.2\)，初始偏差只有 \(e(0) = 0.1\) 度。解是 \(e(t) = 0.1 e^{0.2t}\)：每 3.5 秒翻一倍。20 秒后偏差膨胀到 \(0.1 \times e^4 \approx 5.5\) 度。

如果再加上开窗的冷风 \(w = -1\)，新``平衡点''在 \(e = w/\lambda = -5\) 度——但这个平衡点是山脊而不是山谷：站在上面是平的，稍微偏离就指数滑走。轨迹根本不会停在那里。

这跟散热例子有质的区别。前者是``没人拉住你''——偏差被扰动推着飘。后者是``自己在推自己''——偏差主动放大自身。开环在这两种情形下同样失灵：输入里没有任何关于世界正在发生什么的数据，不管偏差是被推走还是自己跑走，控制律都对此毫无知觉。

这个现象的离散版本出现在 \hyperref[sec:15-Differential-Equations-and-Dynamical-Systems/02-Stability-and-Dynamical-Systems/03]{``离散动力系统、分岔与混沌''}——Logistic 映射对初值的敏感依赖。规律完全相同：指数发散面前，事先算好的轨线只有有限的有效期。你算得再精细、再多位小数，发散的速度都会把它吞掉。

\subsection{直觉的转折：伸手试温}

冲澡时你不需要解任何方程。一边洗一边伸手试，凉了就拧热，烫了就拧冷。翻译成方程，只做一处改动——把测量到的偏差接回输入：

\[
\bar u = -k e, \quad k > 0.
\]

水温低于目标就多加热，高于目标就减热。偏差方程瞬间改写：

\[
\dot e = -(a + k) e + w.
\]

两件事立刻变了。无扰动时 \(e \to 0\) 的收敛率从 \(a\) 跳升到 \(a+k\)——反馈不只是纠偏，还加速纠偏。常值扰动下，稳态偏差从 \(w/a\) 被压缩到

\[
e_{\mathrm{ss}} = \frac{w}{a + k}.
\]

开窗例子中取 \(k = 1\)：\(e_{\mathrm{ss}} = -1 / 1.1 \approx -0.9\) 度。对比开环的 \(-10\) 度——一行 \(\bar u = -ke\) 的改动，残差直接压了一个数量级。

对照本节开头那张框图：开环方案相当于把``测量反馈''那条线剪断，控制器只能读目标 \(r\)；接上这条线，控制器读到的变成误差 \(e\)，而 \(\bar u = -ke\) 就是它接到误差之后做的全部事情。一条信息通路的有无，换来了收敛率与稳态残差两处改变。

模型失配也会在偏差方程中表现为等效扰动——比例反馈同样能把它压小，但不能保证精确归零。要抹掉常值残差，需要下面要讲的积分项。

\subsection{用势能面理解反馈的力学本质}

为什么``把误差接回去''有这种效果？用 \hyperref[sec:15-Differential-Equations-and-Dynamical-Systems/02-Stability-and-Dynamical-Systems/02]{``稳定性与 Lyapunov 函数''} 的语言看得最清晰。构造一个处处非负、在目标处取最小的函数：

\[
V(e) = \frac12 e^2.
\]

沿无扰动的闭环轨迹求导：

\[
\dot V = e \dot e = -(a + k) e^2 < 0 \quad (e \neq 0).
\]

\(V\) 像是一个``高度''或者``势能''——系统沿轨迹滑下这个势能面的斜坡，偏差只能向零沉落。反馈做的事，是把闭环系统改造为一个以目标为汇的梯度流：在任何偏离目标的位置上，动力学内部都产生一个指向零的恢复力。

开环没有这个结构。它的 \(V\) 只被扰动推着走，动力学本身不产生任何指向目标的力。反馈不是在``把计划算得更精细''——它在改变系统的力学结构，从没有恢复力变成有恢复力。

\subsection{PID：三种纠偏力各治一种病}

工程界把纠偏力拆成三项，各管一段时间尺度，各带一种副作用，合起来就是 PID：

\[
u(t) = u^* - k_P e(t) - k_I \int_0^t e(\tau)\, d\tau - k_D \dot e(t).
\]

\textbf{比例管现在。}偏差多大就纠多大。直观、见效快。但有一个结构性的局限：有常值扰动时，稳态必须留下 \(w/(a + k_P)\) 的残差来``养活''纠偏力。偏差归零的瞬间，比例项也为零——驱动中断，扰动又把系统推开。比例控制无法根除静差。一味加大 \(k_P\) 会吃掉稳定裕度，把迟缓的纠正变成剧烈的振荡——像淋浴时冷热水龙头拧得过猛，来回摆荡。

\textbf{积分管历史。}只要偏差没归零，积分项就持续累积，把输入一路推高，直到偏差被彻底抹平。它治的正是比例留下来的静差。代价来自执行器的物理上限：饱和期间误差继续累积，积分项胀到远超实际需要，等误差反号后还要花漫长时间``退账''，造成大幅超调——这叫积分饱和（integrator windup），工程上要对积分项限幅（anti-windup）。积分还引入额外的相位滞后，参数不当本身就能独立地把系统推入振荡。

\textbf{微分管趋势。}它不看偏差有多大，只看偏差变得有多快。水温正快速冲近目标时提前收力，像给系统加阻尼，抑制超调。它的病是噪声：测量噪声的典型成分是高频抖动，而微分恰好放大高频。\(k_D\) 一大，控制信号就被噪声淹没，实际使用中微分项必须配低通滤波器。

三项分工：P 即时纠偏，I 消灭残差，D 预判趋势。调参的实质是在``纠得快''和``稳得住、扛得住噪声''之间做交换——PID 里没有普适的``最优参数''，只有对你当前场景够用的那组权衡。

\subsection{反馈不是万能的}

反馈能在有扰动时做出开环做不到的事。但它的能力边界同样需要用同样诚实的眼光审视。

反馈只对``已经发生的偏差''行动。测量需要时间，纠偏就有延迟。延迟配上过大的增益，把一个稳定系统逼成振荡——淋浴时水温在烫和凉之间来回摆荡，就是延迟加高增益的典型症状。开头那张框图里的反馈线是一条双向通道：偏差信息沿它回到比较点，测量噪声也沿它回到比较点，传感器的颤抖于是变成加热器的颤抖。

更根本的一条：反馈默认``目标已知、偏差可测''。目标为什么是 60 度？如果同时要求省电、升温快、波动小——三个目标彼此冲突，PID 的三个旋钮不是为多目标优化设计的。它稳得住 60 度，但不回答这 60 度本身是否值得追求。

实践中更常见的做法是开环与闭环混合：用模型算一条大方向上的前馈计划，再用反馈清扫剩下的残差。模型管大方向，反馈管小修正，两者分工而非互斥。

把这几条边界并排看，会发现它们指向同一句话：控制结构必须与扰动类型匹配。比例项需要非零偏差来维持抵消常值扰动的力量，因此必留静差；积分项能对常值失配累积出固定补偿，却跟不上任意时变的扰动，还要额外付出相位滞后。这条原则会一路贯穿后面的最优控制设计。

承认这些之后，还剩一个 PID 回答不了的问题：什么叫``好''的控制？稳住只是及格。省燃料、少费时间、少超调，这些追求得先写成一个目标函数，优化才有登场的余地。下一节从这里开始。
